% Section: modal-validity

\documentclass[../../../include/open-logic-section]{subfiles}

\begin{document}

\olfileid[fr]{nml}{syn}{val}

\olsection{Validité}

\begin{explain}
  Les !!{formula}s vraies dans tous les modèles, c'est-à-dire en
  chaque monde de chaque modèle, présentent un intérêt particulier.
  Elles expriment des propositions modales vraies quelle que soit
  l'interprétation de $\Box$ et de $\Diamond$, pourvu que cette
  interprétation soit « normale » : elle est déterminée par une
  relation d'accessibilité entre mondes possibles. Nous appelons
  ces !!{formula}s \emph{valides}. Par exemple,
  $\Box(p \land q) \lif \Box p$ est valide. En revanche, certaines
  !!{formula}s que l'interprétation aléthique de $\Box$ pourrait
  laisser croire valides, comme $\Box p \lif p$, ne le sont pas.
  L'intérêt des modèles relationnels tient notamment au fait que
  différents types de relations d'accessibilité permettent de
  représenter différentes interprétations de $\Box$ et de
  $\Diamond$. Cela suggère de définir la validité non seulement
  relativement à \emph{tous} les modèles, mais aussi relativement
  à tous les modèles \emph{d'un certain type}. Il apparaîtra, par
  exemple, que $\Box p \lif p$ est vraie dans tous les modèles où
  chaque monde est accessible depuis lui-même, autrement dit où
  $R$ est réflexive. Une définition relative aux classes de
  modèles permet de l'énoncer succinctement :
  $\Box p \lif p$ est valide dans la classe des modèles réflexifs.
\end{explain}

\begin{defn}
  !!^a{formula}~$!A$ est \emph{valide} dans une classe
  $\mClass{C}$ de modèles si elle est vraie dans chaque modèle
  de~$\mClass{C}$ (c'est-à-dire en chaque monde de chaque modèle
  de~$\mClass{C}$). Si $!A$ est valide dans~$\mClass{C}$, nous
  écrivons $\mClass{C} \Entails !A$; nous écrivons $\Entails !A$
  si $!A$ est valide dans la classe de \emph{tous} les modèles.
\end{defn}

\begin{prop}\ollabel{prop:subset-class}
  Si $!A$ est valide dans $\mClass{C}$, elle est aussi valide dans
  toute classe $\mClass{C}' \subseteq \mClass{C}$.
\end{prop}

\begin{prop}\ollabel{prop:Nec-rule}
  Si $!A$ est valide, $\Box!A$ l'est également.
\end{prop}

\begin{proof}
  Supposons $\Entails !A$. Pour montrer $\Entails \Box!A$, soient
  $\mModel{M} = \tuple{W, R, V}$ un modèle et $w \in W$. Si $Rww'$,
  alors $\mSat{M}{!A}[w']$, puisque $!A$ est valide. Ainsi,
  $\mSat{M}{\Box!A}[w]$. Comme $\mModel{M}$ et $w$ étaient
  arbitraires, $\Entails \Box!A$.
\end{proof}

\begin{prob}
  Montrez que les !!{formula}s suivantes sont valides :
  \begin{enumerate}
  \item $\Box p \lif \Box (q \lif p)$;
  \item $\Box \lnot \lfalse$;
  \item $\Box p \lif (\Box q \lif \Box p)$.
  \end{enumerate}
\end{prob}

\begin{prob}
  Montrez que $!A \lif \Box!A$ est valide dans la classe
  $\mClass{C}$ des modèles $\mModel{M} = \tuple{W, R, V}$ tels
  que $W = \{w\}$. Montrez de même que $!B \lif \Box !A$ et
  $\Diamond !A \lif !B$ sont valides dans la classe des modèles
  $\mModel{M} = \tuple{W, R, V}$ tels que $R = \emptyset$.
\end{prob}

\end{document}
